
%%%%%%%%%%%%%%%%%%%%%%%%%%%%%%%%%%%%%%%%
%           main body
%%%%%%%%%%%%%%%%%%%%%%%%%%%%%%%%%%%%%%%%


%-------------------------------------
%	TITLE PAGE
%-------------------------------------


\title[Calculus]{\LARGE \bfseries 第一部分\quad 微积分} % The short title appears at the bottom of every slide, the full title is only on the title page
\author[Exercises] % Your institution as it will appear on the bottom of every slide, may be shorthand to save space
{\Large \bfseries 习题}
% \author{Singular Value Decomposition} % Your name
% \date{\today} % Date, can be changed to a custom date
\date{}

%------------------------------------------------

\begin{document}

%------------------------------------------------

%------------------------------------------------

\begin{frame}

\titlepage % Print the title page as the first slide

\end{frame}

%------------------------------------------------

% \begin{frame}
% \frametitle{内容提要} % Table of contents slide, comment this block out to remove it
% \tableofcontents % Throughout your presentation, if you choose to use \section{} and \subsection{} commands, these will automatically be printed on this slide as an overview of your presentation
% \end{frame}

%---------------------------------------
%	PRESENTATION SLIDES
%---------------------------------------

%------------------------------------------------

\begin{frame}

{\color{red} \scriptsize 带$^*$号的题目可以不做}

\vspace{1em}

1. 求下列极限
\begin{itemize}
\item[(1)] $\lim\limits_{x\to-\infty} \dfrac{1}{1+e^{-x}}$
\item[(2)] $\lim\limits_{x\to+\infty} \dfrac{e^x-e^{-x}}{e^x+e^{-x}}$
\item[(3)] $\lim\limits_{x\to 0} \dfrac{\sqrt{1+\tan x} - \sqrt{1-\tan x}}{1-e^{-x}}$
\item[(4)] $\lim\limits_{x\to +\infty} x^{1/x}$
\item[(5)] $\lim\limits_{x\to +\infty} x\left( \frac{1}{e} - \left(\frac{x}{x+1}\right)^x \right)$
\item[(6)] $\lim\limits_{x\to 1-} e^{-\frac{1}{\sqrt{1-x}}}$
\item[(7)] $\lim\limits_{n\to \infty} \dfrac{n!}{n^n}$
\end{itemize}

\end{frame}

%------------------------------------------------

\begin{frame}

2$^*$. 已知$e = \lim\limits_{n\to\infty} \left(1+\frac{1}{n}\right)^n$. 令
\begin{align*}
e_n & = \left(1+\frac{1}{n}\right)^n \\
s_n & = \sum\limits_{k=0}^n \frac{1}{k!} = 1 + \frac{1}{1!} + \frac{1}{2!} + \cdots + \frac{1}{n!} \\
s_{n,k} & = 1 + 1 + \frac{1}{2!}\left( 1-\frac{1}{n} \right) + \cdots + \frac{1}{k!}\left( 1-\frac{1}{n} \right)\cdots \left( 1-\frac{k-1}{n} \right)
\end{align*}
\begin{itemize}
\item[(1)] 写出$e_n$的二项展开。设$n\ge k$, 比较$e_n,s_n,s_{n,k}$三者的大小
\item[(2)] 证明$s_k = \lim\limits_{n\to\infty} s_{n,k}$, 并由此证明$s_k\le e$对任意$k$成立
\item[(3)] 证明$e=\lim\limits_{n\to\infty} s_n$
\end{itemize}

\end{frame}

%------------------------------------------------

\begin{frame}

3. 求下列函数$f$的导数$f'$
\begin{itemize}
\item[(1)] $f(x) = \dfrac{1}{1+e^{-x}}$
\item[(2)] $f(x) = \ln\ln x$
\item[(3)] $f(x) = x^x$
\end{itemize}

\vspace{1em}

4. 若$y = \sin t^2, x = (1 + \sqrt{t})e^t$, 求$\frac{\mathrm{d}y}{\mathrm{d}x}$

\vspace{1em}

5. 求椭圆$\frac{x^2}{a^2} + \frac{y^2}{b^2} = 1$在点$(x_0,y_0)$处的切线方程

\vspace{1em}

6. 设$f(x) = \arctan x$, 求$f^{(n)}(0)$

\vspace{1em}

7. 设$y = e^t, x = \ln t$, 求$\frac{\mathrm{d}^2y}{\mathrm{d}x^2}$

\end{frame}

%------------------------------------------------

\begin{frame}

8. 设$a,b\in\mathbb{R}$为常数，求函数$f(x) = a \ln x + b \ln (1-x)$的极值

\vspace{1em}

9. 设$(a,b)$为曲线$x^4+y^4-x^2y-xy^2 = 0$上的动点，求$b$能取到的最大值与最小值

\vspace{1em}

10. 试写出函数$f(x) = \ln\cos x$在$x=0$处的4阶泰勒多项式

\vspace{1em}

11. 求函数$f(x) = \ln(1+e^{-x})$的二阶导数$f''(x)$, 并由此判断该函数的凹凸性

\vspace{1em}

12. 考虑函数$\displaystyle \Gamma(x) = \int_0^{+\infty} \dfrac{t^{x-1}}{e^t} dt, \; x\in\mathbb{R}$
\begin{itemize}
\item[(1)] 证明$\Gamma(x+1) = x\Gamma(x)$
\item[(2)] 求$\Gamma(1)$
\end{itemize}

\end{frame}

%------------------------------------------------

\begin{frame}

13. 计算以下积分
\begin{itemize}
\item[(1)] $\displaystyle \int_0^1 \dfrac{x}{1+\sqrt{1+x}} \mathrm{d}x$
\item[(2)] $\displaystyle \int_0^1 x^3e^{-x^2} \mathrm{d}x$
\item[(3)] $\displaystyle \int_0^{+\infty} \dfrac{1}{1+e^x} \mathrm{d}x$
\item[(4)] $\displaystyle \int_0^1 \sqrt{\dfrac{x}{1-x}} \mathrm{d}x$
\item[($5^*$)] $\displaystyle \int_0^{+\infty} e^{-x^2} \mathrm{d}x$
\end{itemize}

\end{frame}

%------------------------------------------------

\begin{frame}

14. 求偏导
\begin{itemize}
\item[(1)] $f(x,y)=y^{\ln x}$, 求$\dfrac{\partial^2 f}{\partial x \partial y}$
\item[(2)] 设$z = \dfrac{1}{1+e^{x-2y}}, x = \cos t, y = \sqrt{1+t^2}$, 求$\dfrac{dz}{dt}$
\item[(3)] 设$E(w,b) = \sum\limits_{i=1}^m (y_i-wx_i-b)^2,$ 求$\dfrac{\partial E}{\partial w}, \dfrac{\partial E}{\partial b}$
\item[(4)] 设$y(w_1,\ldots,w_n,b) = \dfrac{1}{1+e^{-\sum\limits_{i=1}^n w_ix_i - b}}$, 求$\dfrac{\partial y}{\partial w_i}$
\end{itemize}

\end{frame}

%------------------------------------------------

\begin{frame}

15. 设$f(x) = \ln \left(e^{-x} + \dfrac{x^2}{y} \right)$
\begin{itemize}
\item[(1)] 求$f(x)$的梯度向量
\item[(2)] 令$\mathbf{v} = (1,1)$, 求$\left. \dfrac{\partial f}{\partial \mathbf{v}} \right|_{(-1,-1)}$
\end{itemize}

\vspace{1em}

16. 在神经网络的反向传播算法中（见讲义\S 9），取损失函数为$L(y,\widehat{y}) = \frac{1}{2}(y-\widehat{y})^2$, 激励函数为$\sigma(x) = \dfrac{1}{1+e^{-x}}$, 学习率取作$\eta$, 求权重$w_{ij}$, $v_{ij}$在每一轮训练中更新的增量的具体表达式

\end{frame}

% %------------------------------------------------
% % closing page
% \begin{frame}
% \HUGE{\centerline{\bfseries {\color{gray} The} {\color{zkblue} End}}}

% \pause

% \vspace{1.2em}

% \Large{\centerline{\bfseries {\color{gray}Thank} \quad {\color{zkblue} You}}}

% \end{frame}

%----------------------------------------------------------------------------------------

\end{document}
